This draft is heavily work in progess.
Here is a quick rundown of the strategy (not everything has been tex'ed yet, but all statements have been worked out) for proving Serre duality for vector bundles.
This is most likely subject to change as the ideas settle.

We first prove Serre duality for vector bundles for $\bP^n$.
This generality is not very difficult to see and essentially follows the classical argument aside form a couple of minor adjustments.
Then we generalize to a smooth projective scheme $X$ of pure dimension.
For this, we also follow a known route that is certainly not novel.
However, even though in the course of the argument we pass by the same major theorem that a classical mathematician would, the proof and methods we use to obtain these results are quite different.
Here are the main results used for a general $X$:

We show in \Cref{merely-finite-map}, that any such smooth projective scheme $X$ of pure dimension merely admits a finite map $f : X \to \bP^{\dim X}$.
To prove this, we will, with \Cref{weak-bertini}, prove a version of a folklore result by Bertini, establishing that the set of hyperplanes $H$ intersecting $X$ such that $H \cap X$ is smooth of dimension $\dim X - 1$ is open and not empty.
We then employ (and develope) the theory of \emph{étale-overt} schemes to upgrade the double negated existence of such a hyperplane to an étale-local one.
This stronger existence turns out to be enough to merely get the datum required to construct the desired finite map.
It is an interesting observation, that the étale-local existence of the finite map $f$, while being substantially easier to show, is also surprisingly useful: Lots of interesting propositions in the vicinity of Serre duality are in fact already étale sheaves.

We will also establish, in \Cref{finite-map-of-equidimensional-smooth-schemes-preserves-vectorbundles}, that $f_\ast$ preserves vector bundles and, when retricted to these, admits a right adjoint $f^!$.
For the remaining argument we then rejoint the well travelled road of the classical argument and mostly use formal reasoning to deduce the general case of Serre duality with $f^!\omega_{\bP^n}$ as the dualizing sheaf.
\bigskip

This draft is currently not meant to be the shortest possible path to Serre duality and we instead take our time to develope some other results in which are interesting in their own right.
\bigskip

In this treatment we use the notion of a \emph{finitely adequate} $R$-module, which we hope will serve as a satsfactory replacement of the faulty category of coherent modules in our setting.
It is being developed in parallel in a different draft \href{https://felix-cherubini.de/adequate-modules.pdf}{here}, together with Dan Christensen, Hugo Moeneclay and Thomas Thorbjørnsen.
